\documentclass[../main]{subfiles}

% Appendix A: a plain appendix (no opening page, no epigraph). Labels use
% the app: prefix. Refer to an appendix as a whole ("Appendix~\ref{...}")
% from the chapters rather than to its equations or figures.

\begin{document}
% In a standalone build, switch to appendix numbering (A, B, ...).
\ifSubfilesClassLoaded{\appendix}{}
\chapter{An example appendix}
\label{app:example}

This appendix collects technical details that Chapter~\ref{ch:ch1} leaves out.
Acronyms count per appendix as they do per chapter: a single use is spelled out with \verb|\acrlong|, as the \acrlong{cpr} here.
\lipsum[1][1-4]
In Section~\ref{app:derivation} we derive an expression, and in Section~\ref{app:limits} we discuss its limits.

\section{A derivation}
\label{app:derivation}
\lipsum[2][1-4]
Integrals keep a thin space before the differential, as in
\begin{equation}
    \label{eq:app_integral}
    N = \int_0^R dr\,r\rho(r) = \int d\omega\, f(\omega)\nu(\omega),
\end{equation}
which is the only place a thin space appears in math.
\lipsum[3][1-5]

\section{Limiting cases}
\label{app:limits}
\lipsum[4][1-6]

\lipsum[5][1-6]

\ifSubfilesClassLoaded{%
\newrefcontext[labelprefix=]
\printbibliography
}{}
\end{document}
